% !TEX TS-program = XeLaTeX
% use the following command:
% all document files must be coded in UTF-8
\documentclass[portuguese]{textolivre}
% build HTML with: make4ht -e build.lua -c textolivre.cfg -x -u article "fn-in,svg,pic-align"

\journalname{Texto Livre}
\thevolume{19}
%\thenumber{1} % old template
\theyear{2026}
\receiveddate{\DTMdisplaydate{2025}{12}{29}{-1}} % YYYY MM DD
\accepteddate{\DTMdisplaydate{2026}{3}{18}{-1}}
\publisheddate{\DTMdisplaydate{2026}{8}{28}{-1}}
\corrauthor{Maria Aparecida Araujo Lima}
\articledoi{10.1590/1983-3652.2026.63671}
%\articleid{NNNN} % if the article ID is not the last 5 numbers of its DOI, provide it using \articleid{} commmand 
% list of available sesscions in the journal: articles, dossier, reports, essays, reviews, interviews, editorial
\articlesessionname{articles}
\runningauthor{Araujo Lima, Muñoz-Repiso e Meirinhos} 
%\editorname{Leonardo Araújo} % old template
\sectioneditorname{Daniervelin Pereira~\orcid{0000-0003-1861-3609}}
\layouteditorname{Rayany Chaves~\orcid{0009-0007-0028-484X}}

\title{Tradutores da língua de sinais e possibilidade de uso por professores e estudantes do ensino superior}
\othertitle{Sign Language Translators and the Possibility of Use by Higher Education Teachers and Students}
% if there is a third language title, add here:
%\othertitle{Artikelvorlage zur Einreichung beim Texto Livre Journal}

\author[1,2]{Maria Aparecida Araujo Lima~\orcid{0000-0001-7187-2373}\thanks{Email: \href{mailto:cidaaraujo.uab@gmail.com}{cidaaraujo.uab@gmail.com}}}
\author[1]{Ana Garcia-Valcarcel Muñoz-Repiso~\orcid{0000-0003-0463-0192}\thanks{Email: \href{mailto:anagv@usal.es}{anagv@usal.es}}}
\author[2]{Manuel Meirinhos~\orcid{0000-0003-1756-709X}\thanks{Email: \href{mailto:meirinhos@ipb.pt}{meirinhos@ipb.pt}}}
\affil[1]{Universidade de Salamanca, Faculdade de
Educação, Departamento de Didática, Organização e Métodos de Pesquisa, Paseo de Canalejas, Salamanca, Espanha.}
\affil[2]{Universidade Politécnica de Bragança, Centro de Investigação Transdisciplinar em Educação e Desenvolvimento (CITeD), Bragança, Portugal.}

\addbibresource{article.bib}
% use biber instead of bibtex
% $ biber article

% used to create dummy text for the template file
\definecolor{dark-gray}{gray}{0.35} % color used to display dummy texts
\usepackage{lipsum}
\SetLipsumParListSurrounders{\colorlet{oldcolor}{.}\color{dark-gray}}{\color{oldcolor}}

% used here only to provide the XeLaTeX and BibTeX logos
\usepackage{hologo}

% if you use multirows in a table, include the multirow package
\usepackage{multirow}

% provides sidewaysfigure environment
\usepackage{rotating}

% CUSTOM EPIGRAPH - BEGIN 
%%% https://tex.stackexchange.com/questions/193178/specific-epigraph-style
\usepackage{epigraph}
\renewcommand\textflush{flushright}
\makeatletter
\newlength\epitextskip
\pretocmd{\@epitext}{\em}{}{}
\apptocmd{\@epitext}{\em}{}{}
\patchcmd{\epigraph}{\@epitext{#1}\\}{\@epitext{#1}\\[\epitextskip]}{}{}
\makeatother
\setlength\epigraphrule{0pt}
\setlength\epitextskip{0.5ex}
\setlength\epigraphwidth{.7\textwidth}
% CUSTOM EPIGRAPH - END

% to use IPA symbols in unicode add
%\usepackage{fontspec}
%\newfontfamily\ipafont{CMU Serif}
%\newcommand{\ipa}[1]{{\ipafont #1}}
% and in the text you may use the \ipa{...} command passing the symbols in unicode

% LANGUAGE - BEGIN
% ARABIC
% for languages that use special fonts, you must provide the typeface that will be used
% \setotherlanguage{arabic}
% \newfontfamily\arabicfont[Script=Arabic]{Amiri}
% \newfontfamily\arabicfontsf[Script=Arabic]{Amiri}
% \newfontfamily\arabicfonttt[Script=Arabic]{Amiri}
%
% in the article, to add arabic text use: \textlang{arabic}{ ... }
%
% RUSSIAN
% for russian text we also need to define fonts with support for Cyrillic script
% \usepackage{fontspec}
% \setotherlanguage{russian}
% \newfontfamily\cyrillicfont{Times New Roman}
% \newfontfamily\cyrillicfontsf{Times New Roman}[Script=Cyrillic]
% \newfontfamily\cyrillicfonttt{Times New Roman}[Script=Cyrillic]
%
% in the text use \begin{russian} ... \end{russian}
% LANGUAGE - END

% EMOJIS - BEGIN
% to use emoticons in your manuscript
% https://stackoverflow.com/questions/190145/how-to-insert-emoticons-in-latex/57076064
% using font Symbola, which has full support
% the font may be downloaded at:
% https://dn-works.com/ufas/
% add to preamble:
% \newfontfamily\Symbola{Symbola}
% in the text use:
% {\Symbola }
% EMOJIS - END

% LABEL REFERENCE TO DESCRIPTIVE LIST - BEGIN
% reference itens in a descriptive list using their labels instead of numbers
% insert the code below in the preambule:
%\makeatletter
%\let\orgdescriptionlabel\descriptionlabel
%\renewcommand*{\descriptionlabel}[1]{%
%  \let\orglabel\label
%  \let\label\@gobble
%  \phantomsection
%  \edef\@currentlabel{#1\unskip}%
%  \let\label\orglabel
%  \orgdescriptionlabel{#1}%
%}
%\makeatother
%
% in your document, use as illustraded here:
%\begin{description}
%  \item[first\label{itm1}] this is only an example;
%  % ...  add more items
%\end{description}
% LABEL REFERENCE TO DESCRIPTIVE LIST - END


% add line numbers for submission
%\usepackage{lineno}
%\linenumbers

\begin{document}
\maketitle

\begin{polyabstract}
\begin{abstract}
Este estudo teve como objetivo analisar e comparar sete tradutores automáticos de língua de sinais baseados em IA: \textit{Dicionário da Língua Brasileira de Sinais}, \textit{Spreadthesign}, \textit{Glove-Based Systems}, \textit{Hand Talk}, \textit{VLibras}, \textit{Sign-Speak} e CESAR + Lenovo. A metodologia utilizada foi qualitativa, exploratória, realizada em duas etapas: (1) levantamento bibliográfico e análise comparativa com critérios de clareza e (2) coerência e fidelidade semântica. Os tradutores foram selecionados por três critérios: (1) recurso dicionarizado, (2) tradução automática online e (3) aplicabilidade institucional, educacional e corporativa. Os resultados indicaram que o \textit{Dicionário da Língua Brasileira de Sinais} e \textit{Spreadthesign} não atuam como tradutores, mas ampliam a acessibilidade de pessoas surdas. \textit{Glove-Based Systems} traduz apenas da língua de sinais para língua oral/escrita; \textit{Hand Talk} e \textit{VLibras} traduzem somente do português para sinais. \textit{Sign-Speak} e CESAR + Lenovo permitem tradução bidirecional do inglês ou português para ASL ou Libras, (e vice-versa), com reconhecimento por câmera. Conclui-se que não há tradutor que atenda plenamente a bidirecionalidade e a precisão, indicando uma necessidade de avanços em IA para acessibilidade comunicacional. A tradução automática de línguas de sinais no Brasil enfrenta um desafio importante: a diversidade regional da Libras, que dificulta padronização dos sistemas. Além disso, uso de IA em educação deve ser planejado com cuidado, considerando limitações e integração a práticas pedagógicas inclusivas, com participação ativa de usuários fluentes. Destaca-se ainda que esses sistemas variam quanto ao desempenho, interface e usabilidade, influenciando a experiência do usuário. Observa-se também que limitações tecnológicas atuais impactam a qualidade das traduções, especialmente em contextos complexos e interações em tempo real, reforçando a importância de pesquisas contínuas na área para maior inclusão.

\keywords{Dicionários de tradução de língua de sinais \sep Luvas de tradução de língua de sinais \sep Tradutores de sinais unidirecional e bidirecional \sep Educação inclusiva \sep Acessibilidade}
\end{abstract}

\begin{english}
\begin{abstract}
The aim of this study was to analyze and compare seven AI-based sign language machine translators: \textit{Brazilian Sign Language Dictionary}, \textit{Spreadthesign}, \textit{Glove-Based Systems}, \textit{Hand Talk}, \textit{VLibras}, \textit{Sign-Speak} and CESAR + Lenovo. The methodology used was qualitative and exploratory, consisting of two stages: (1) a literature review and a comparative analysis using criteria of clarity and (2) coherence and semantic fidelity. The translators were selected based on three criteria: (1) dictionary-based resources, (2) online machine translation and (3) institutional, educational and corporate applicability. The results indicated that the \textit{Brazilian Sign Language Dictionary} and \textit{Spreadthesign} do not function as translators but enhance accessibility for deaf people. \textit{Glove-Based Systems} translates only from sign language to spoken or written language; \textit{Hand Talk} and \textit{VLibras} translate only from Portuguese to sign language. \textit{Sign-Speak} and CESAR + Lenovo enable bidirectional translation from English or Portuguese to ASL or Libras (and vice versa), using camera recognition. It can be concluded that no translator fully meets the requirements for bidirectionality and accuracy, indicating a need for advances in AI for communication accessibility. Machine translation of sign languages in Brazil faces a significant challenge: the regional diversity of Libras, which hinders the standardization of systems. Furthermore, the use of AI in education must be carefully planned, considering its limitations and integration into inclusive teaching practices, with active participation of fluent users. These systems also vary in performance, interface, and usability, influencing user experience. Current technological limitations affect translation quality, especially in complex contexts and real-time interactions, reinforcing the importance of ongoing research to promote greater inclusion.

\keywords{Sign language translation dictionaries \sep Sign language translation gloves \sep Unidirectional and bidirectional sign translators \sep Inclusive education \sep Accessibility}
\end{abstract}
\end{english}
% if there is another abstract, insert it here using the same scheme
\end{polyabstract}

\section{Introdução}
A língua de sinais constitui um dos pilares da identidade cultural, linguística e social das pessoas surdas, desempenhando um papel essencial na comunicação e na inclusão em diferentes contextos educativos, profissionais ou sociais. Ela é uma língua completa, com gramática e estrutura própria, que reflete a diversidade cultural e regional dos povos que fazem uso dela.

No Brasil, por exemplo, a Língua Brasileira de Sinais (Libras) é reconhecida legalmente e representa um marco no avanço dos direitos linguísticos da comunidade surda. Contudo, persistem barreiras significativas à comunicação entre surdos e ouvintes, especialmente em ambientes educacionais, corporativos e institucionais, o que reforça a urgência de soluções tecnológicas acessíveis e sensíveis à diversidade linguística.

Para que não haja obscuridade por parte do leitor, esclarecemos a diferenciação entre pessoas surdas e pessoas com deficiência auditiva, com base no Decreto n.º 5.626/2005, que estabelece a diferença em seu Artigo n.º 2: 

\begin{quote}
    Art. 2º Para os fins deste Decreto, considera-se pessoa surda aquela que, por ter perda auditiva, compreende e interage com o mundo por meio de experiências visuais, manifestando sua cultura principalmente pelo uso da Língua Brasileira de Sinais - Libras.
    
    Parágrafo único. Considera-se deficiência auditiva a perda bilateral, parcial ou total, de quarenta e um decibéis (dB) ou mais, aferida por audiograma nas freqüências de 500Hz, 1.000Hz, 2.000Hz e 3.000Hz \cite{Brasil2005Decreto5626}.
\end{quote}

%Precisa ter espaçamento entre as frases da citação, para ficarem como o autor colocou no artigo

Entendemos que a interação entre surdos e ouvintes requer compreensão das particularidades linguísticas da língua de sinais formal e informal usada por pessoas com deficiência auditiva, que impacta diretamente a participação dos alunos surdos no ensino superior, exigindo adaptações, recursos acessíveis e a inserção prática da temática da inclusão nos cursos universitários. Diante disso, formulamos a seguinte questão: os tradutores da língua de sinais podem ser usados por professores e estudantes do ensino superior como recurso inclusivo para facilitar a comunicação entre pessoas surdas e ouvintes?

Na universidade, o uso da escrita acadêmica estimula o avanço da gramaticalização da Libras, impulsionado por produções acadêmicas de surdos, intérpretes e tutores online. Os tradutores automáticos de sinais contribuem para o fortalecimento das relações entre surdos e ouvintes, favorecendo a disseminação e aprimoramento dessas tecnologias conforme as especificidades de cada área do conhecimento.

De acordo com \textcite{QuadrosKarnopp2004,Skliar1998}, o uso da Libras em contextos acadêmicos contribui significativamente para o desenvolvimento de sua estrutura linguística. Quando inserida em ambientes formais, como a universidade, a língua de sinais passa a exigir maior complexidade, o que favorece a ampliação do vocabulário, a padronização de usos e o fortalecimento de sua gramática. O uso acadêmico da Libras, por exemplo, exige a expressão de conceitos abstratos e específicos, promovendo seu aprimoramento. Portanto, a universidade desempenha um papel fundamental na evolução e fortalecimento da língua de sinais.

Em uma situação acadêmica, para um estudante surdo apresentar um trabalho sobre “educação inclusiva”, ele vai precisar explicar conceitos abstratos, organizar argumentos, usar termos específicos (ex.: “metodologia”, “hipótese”, “análise”). Nesse contexto, a Libras evolui pela necessidade de produções formais, como trabalhos acadêmicos, artigos e apresentações, que demandam clareza, organização e precisão. Esse processo estimula a gramaticalização da língua de sinais.

Além disso, essa evolução é impulsionada pelos sujeitos, como pessoas surdas que produzem conhecimento diretamente em Libras, intérpretes que fazem a tradução entre a Libras e o português, e vice-versa, e tutores que ajudam na circulação do conteúdo (presencial ou online).Todos contribuem para a criação de novos sinais, ampliação do vocabulário, adaptação de termos e consolidação de formas mais formais da língua de sinais.

Segundo \textcite{HolmesPorayskaPomsta2023}, a Inteligência Artificial pode contribuir para a inclusão e a acessibilidade educacional por meio de sistemas capazes de reconhecer padrões, processar informações linguísticas e adaptar-se às necessidades dos usuários. Nesse contexto, a IA tem potencial para aprimorar tecnologias de tradução da língua de sinais, favorecendo a comunicação entre pessoas surdas e ouvintes no ambiente educacional.

Ao utilizar técnicas de reconhecimento de padrões e processamento de linguagem, esses sistemas podem tornar a tradução mais ágil e acessível. No entanto, a precisão dessas ferramentas ainda depende da qualidade dos dados utilizados em seu desenvolvimento e das características do contexto de uso. Por essa razão, os resultados gerados por sistemas de IA devem ser avaliados de forma crítica, especialmente em situações que envolvam interpretações complexas da língua de sinais, a fim de garantir uma comunicação adequada e inclusiva.

Tradutores automáticos de língua de sinais, dicionários digitais e sistemas de reconhecimento de sinais vêm sendo desenvolvidos com o objetivo de facilitar a interação das pessoas surdas com as pessoas ouvintes, ampliando o acesso à informação, à educação e à participação social. Esses recursos podem contribuir para que estudantes surdos participem de forma mais autônoma do processo de ensino-aprendizagem e estejam mais integrados em contextos de trabalho e, além disso, para que professores, tutores e intérpretes disponham de novas alternativas didáticas e comunicacionais adaptadas às suas realidades.

Os tradutores automáticos de língua de sinais contribuem para fortalecer as relações entre pessoas surdas e ouvintes, porque reduzem as barreiras de comunicação, tendo em vista que eles traduzem fala ou texto para a língua de sinais. Alguns tradutores possibilitam a tradução da língua de sinais para fala ou texto. Os tradutores também promovem a inclusão social, considerando que funcionam como mediadores comunicacionais, ajudando na inclusão em contextos sociais e educacionais. 

Eles favorecem a participação em aulas e eventos, acesso a serviços públicos e maior interação social e promovem maior autonomia para pessoas surdas, porque são acessíveis e fáceis de usar, o que aumenta a independência dos usuários surdos no dia a dia, bem como a aproximação entre culturas (surda e ouvinte), visto que essas tecnologias ajudam a criar mais empatia, maior compreensão cultural e relações mais naturais entre surdos e ouvintes.

No entanto, considera-se importante mencionar que o uso de avatares em tradutores de língua de sinais enfrenta desafios significativos, especialmente na reprodução das expressões não manuais. Estas expressões, como movimentos faciais e corporais, são essenciais para transmitir emoções, intenções e nuances gramaticais. Os avatares ainda apresentam limitações na precisão desses elementos, comprometendo a compreensão da mensagem; sendo assim, melhorar sua expressividade é crucial para garantir acessibilidade e comunicação eficaz para pessoas com deficiência auditiva.

As técnicas de Inteligência Artificial exigem vasto conhecimento e mecanismos simbólicos para permitir que computadores interpretem, analisem e tomem decisões \cite{Santos2018}. No contexto da acessibilidade, especialmente para pessoas surdas, os tradutores de sinais sejam humanos ou digitais, desempenham papel essencial na comunicação e na inclusão, sobretudo na educação.

Os tradutores automáticos de línguas de sinais, embora representem um avanço importante na promoção da acessibilidade, ainda enfrentam diversas limitações, que exigem um uso crítico por parte dos usuários. Diferentemente das línguas orais, as línguas de sinais possuem um forte componente visual e espacial, envolvendo expressões faciais, movimentos corporais e configurações de mão que são difíceis de captar e traduzir com precisão por sistemas automatizados. 

Além disso, variações regionais e contextuais podem comprometer ainda mais a qualidade da tradução, uma vez que estudos como o de \textcite{XavierBarbosa2014} demonstram a existência de variação sistemática, tanto entre sinalizantes quanto no próprio uso individual na produção dos sinais da Libras, o que representa um desafio adicional para sistemas de tradução automática treinados em conjuntos padronizados de dados. Como destaca \textcite{Rodrigues2018}, a interpretação em línguas de sinais é um processo complexo que vai além da simples correspondência de sinais, envolvendo aspectos operacionais e cognitivos distintos dos da tradução propriamente dita, e exigindo do profissional o desenvolvimento do que o autor chama de Competência Tradutória. Nesse sentido, os tradutores automáticos devem ser entendidos como ferramentas de apoio, úteis para facilitar a comunicação inicial, mas insuficientes para garantir fidelidade e nuance, tornando indispensável uma abordagem crítica e, sempre que possível, a mediação de profissionais qualificados.

Nesse sentido, os tradutores automáticos devem ser entendidos como ferramentas de apoio, úteis para facilitar a comunicação inicial, mas insuficientes para garantir fidelidade e nuance, tornando indispensável uma abordagem crítica e, sempre que possível, a mediação de profissionais qualificados. 

Isso significa não confiar cegamente na tecnologia, avaliando sempre se a tradução faz sentido, especialmente em contextos sensíveis.

O tradutor de sinais interpreta uma frase e pode não ser fiel ao que foi dito pelo emissor, pois ainda não está totalmente preparado para traduzir com fidelidade. A seguir, na \Cref{tab1}, estão alguns exemplos para complementar a explicação:

\begin{table}[htbp]
    \centering
    \caption{Exemplo de dificuldades do tradutor automático de sinais.}
    \label{tab1}
    \begin{tabularx}{\textwidth}{>{\hsize=1.4\hsize\linewidth=\hsize}X>{\hsize=0.6\hsize\linewidth=\hsize}X}
        \toprule
        Contexto & Risco \\ \midrule
        \arrayrulecolor[gray]{.7}
        \emph{Contexto médico} \newline
        Frase: ``Estou com uma dor aguda no peito que irradia para o braço.'' Pode gerar sinais que significam apenas ``dor no peito'' de forma genérica, perdendo a gravidade clínica da expressão ``irradia para o braço'', que pode indicar um problema cardíaco. & 
        Diagnóstico incorreto ou atraso no atendimento. \\ \midrule
        
        \emph{Contexto jurídico} \newline
        Frase: ``O réu agiu sob coação.'' Um tradutor automático pode interpretar como ``foi forçado'', mas sem captar o peso jurídico específico de ``coação'' no Direito. & 
        Interpretação errada de um depoimento ou decisão judicial. \\ \midrule
        
        \emph{Diferenças culturais} \newline
        Frase: ``Ele quebrou o galho para mim.'' Um sistema automático pode traduzir literalmente (algo como ``quebrar um ramo''), sem entender que significa ``ajudar''. & 
        Perda de sentido e confusão na comunicação. \\ \midrule
        
        \emph{Uso educacional} \newline
        Contexto: Um estudante surdo usa um tradutor automático para acompanhar uma aula complexa. O sistema traduz palavras isoladas, mas não consegue estruturar o raciocínio completo nem adaptar explicações. & 
        Compreensão superficial do conteúdo.\\ 
        \arrayrulecolor{black}
        \bottomrule 
        \end{tabularx}
\begin{flushleft}
    \source{Elaborada pelo autor.}
    \end{flushleft}
    \end{table}

Esses exemplos mostram que os tradutores ajudam no básico (comunicação inicial). Entretanto, para garantir que tradução automática ocorra como precisão, considerando o contexto e as nuances cultural e linguística, seriam necessários ajustes para que a IA fizesse o papel de profissionais humanos qualificados em áreas específicas, para que os tradutores automáticos realizassem a tradução real e fidedigna.

A tradução automática online em língua de sinais é uma tecnologia que busca reduzir barreiras de comunicação entre pessoas surdas e ouvintes, permitindo maior interação. Ela utiliza recursos como inteligência artificial, visão computacional e processamento de linguagem para converter, em tempo real (ou quase), línguas faladas/escritas em línguas de sinais, e vice-versa. O sistema envolve o reconhecimento de movimentos (mãos, rosto e corpo), sua interpretação pela IA, adaptação do conteúdo para a outra língua e a geração de saída em texto, voz ou por meio de avatares que reproduzem sinais. Nesse sentido, apesar dos avanços, a IA ainda está em desenvolvimento e não substitui totalmente intérpretes humanos.

Assim, compreender as potencialidades e restrições dos tradutores de sinais é fundamental para garantir uma comunicação eficaz e equitativa entre surdos e ouvintes, promovendo avanços contínuos na acessibilidade e na inclusão educacional e social. Outra questão que deve ser considerada pela IA na tradução para pessoas surdas é o fato de que existem diferentes línguas de sinais em países que partilham a língua portuguesa. Isso revela que essas línguas não derivam diretamente da língua oral dominante, mas das interações sociais e culturais das comunidades surdas locais. A \Cref{tab2} apresenta os países da Comunidade dos Países de Língua Portuguesa (CPLP).

\begin{table}[htbp]
    \centering
    \caption{Lista dos países da CPLP e suas respectivas línguas de sinais.}
    \label{tab2}
    \begin{tabular}{lll}
\toprule
País & Língua de sinais & Sigla \\ \midrule
Angola & Língua Angolana de Sinais & LAS \\
Brasil & Língua Brasileira de Sinais & LIBRAS \\ Moçambique & Língua Moçambicana de Sinais & LMS \\
Portugal & Língua Gestual Portuguesa & LGP \\
Cabo Verde & Língua Gestual Cabo-verdiana (em desenvolvimento) & LGCV \\
Guiné-Bissau & Língua Gestual da Guiné-Bissau (não padronizada) & LGGB \\
São Tomé e Príncipe & Língua Gestual São-Tomense (em desenvolvimento) & LGSTP \\
Timor-Leste & Língua Gestual de Timor-Leste (em desenvolvimento) & LGTL \\
Guiné Equatorial & Língua Gestual da Guiné Equatorial influência espanhola) & LGGE \\ 
\bottomrule
\end{tabular}
\source{\textcite{CPLP,Vanali2016}, adaptada pelos autores.}
\end{table}

Segundo \textcite{LopesGoettert2015}, no âmbito educacional, nos mais variados níveis de ensino, a língua de sinais geralmente é desconhecida pelos ouvintes e os professores ouvintes, em sua maioria, também estão inseridos nesse grupo. Mesmo assim, sejam os professores surdos ou ouvintes, estes necessitarão de elaborar, produzir e dispor material didático inclusivo para o uso dos alunos com deficiência auditiva, desde a educação básica até os cursos de ensino superior. Nem sempre é dada a devida capacitação aos estudantes dos cursos de formação de professores na universidade.

O Brasil foi selecionado como exemplo para motivar a reflexão sobre a importância da inclusão de pessoas surdas no contexto educacional e seus impactos na sociedade, considerando que no país, a Língua Brasileira de Sinais (Libras) é reconhecida como língua oficial, por meio da Lei n.º 10.436. Segundo dados do Ministério da Educação (MEC), no Censo da Educação Superior de 2022 foram registradas 11.313 matrículas de pessoas com deficiência auditiva ou surdas em cursos de graduação\footnote{\url{https://download.inep.gov.br/educacao_superior/censo_superior/documentos/2022/apresentacao_censo_da_educacao_superior_2022.pdf}}

Esse cenário evidencia a necessidade de incorporar recursos didáticos baseados em IA, como tradutores de sinais, com o objetivo de ampliar a disseminação e a compreensão da Libras. Dessa forma, busca-se promover maior inclusão e acessibilidade, visando garantir a participação equitativa da comunidade surda no ensino superior.

O presente artigo tem como público-alvo principal as pessoas surdas e os profissionais envolvidos em sua formação e acompanhamento, incluindo professores, tutores, intérpretes, gestores educacionais e desenvolvedores de tecnologias assistivas. Busca-se fornecer subsídios teóricos, mas essencialmente práticos, que orientem o uso consciente e crítico dos tradutores automáticos e demais ferramentas digitais de acessibilidade na comunicação.

O trabalho está estruturado em cinco partes: a primeira apresenta uma abordagem teórica sobre a relevância da língua de sinais e o papel da IA na inclusão comunicacional; a segunda descreve a metodologia de análise comparativa aplicada; a terceira descreve e apresenta os resultados obtidos; a quarta realiza a análise interpretativa dos resultados a partir da avaliação de sete tradutores de sinais: \textit{Dicionário da Língua Brasileira de Sinais, Spreadthesign, Glove-Based Systems, Hand Talk, VLibras, Sign-Speak} e CESAR+ e, por fim, a quinta seção apresenta as considerações finais e recomendações para pesquisas futuras e políticas de inclusão digital voltadas à comunidade surda.

\section{Metodologia}
Este estudo teve como base a seguinte questão: os tradutores da língua de sinais podem ser usados por professores e estudantes do ensino superior como recurso inclusivo, para facilitar a comunicação entre pessoas surdas e ouvintes? Partindo dessa questão, formulou-se um objetivo geral, voltado para a identificação e análise de tradutores de sinais, que facilitasse a comunicação entre surdos e ouvintes em ambientes educacionais universitários. De forma específica, pretendemos analisar e comparar sete tradutores de língua de sinais, com base em três eixos analíticos principais: (1) recursos de dicionário (tradutor automático); (2) tradução automática unidirecional; e (3) tradução automática bidirecional e sua aplicabilidade em contextos institucionais, educacionais e corporativos.

Os tradutores selecionados foram: \textit{Dicionário da Língua Brasileira de Sinais, Spreadthesign, Glove-Based Systems, Hand Talk, VLibras, Sign-Speak} e CESAR + Lenovo. A escolha dessas ferramentas fundamentou-se na sua disponibilidade pública e/ou relevância tecnológica e científica para os contextos educacional e corporativo, tanto em âmbito nacional (Brasil) quanto internacional.

A investigação adotou uma abordagem qualitativa, de caráter exploratório e descritivo. Essa abordagem foi considerada adequada por permitir a compreensão aprofundada de um fenômeno emergente e ainda pouco estudado, privilegiando a interpretação e descrição detalhada em vez da quantificação \cite{CreswellPoth2018}. Dado tratar-se de um domínio em rápida evolução tecnológica, a perspetiva exploratória possibilitou mapear tendências, identificar limitações e apontar oportunidades de desenvolvimento científico, tecnológico e prático no campo dos tradutores automáticos de língua de sinais.

O estudo fundamentou-se em análise documental e comparativa. O corpus foi composto pelos tradutores automáticos de língua de sinais selecionados, atualmente disponíveis em contextos nacionais e internacionais. A recolha de dados baseou-se exclusivamente em fontes secundárias, incluindo websites institucionais, publicações técnicas, artigos científicos, relatórios e vídeos demonstrativos. Conforme \textcite{Bowen2009}, a análise documental é um procedimento válido na investigação qualitativa, permitindo extrair significados, identificar padrões e compreender representações sociais a partir de materiais textuais e multimodais.

A análise comparativa foi conduzida a partir de um conjunto de critérios previamente definidos, que funcionaram como categorias orientadoras – unidirecionalidade (escrita ou voz-para-sinais) e bidirecionalidade da tradução (voz-para-sinais e sinais-para-voz); fidelidade linguística; clareza semântica; usabilidade; aplicabilidade educativa, para ter uma visão da situação da estrutura e da fidelidade da tradução para usuários surdos. Para clarificar a expressão “fidelidade da tradução para usuários surdos”, esclarecemos que ela refere-se ao grau em que a tradução preserva com precisão o conteúdo, o sentido e as nuances da mensagem original ao ser convertida entre modalidades (por exemplo, de voz para língua de sinais, e vice-versa), considerando as especificidades linguísticas e culturais da comunidade surda.

Além da análise documental, também foram realizados acessos nos aplicativos para verificar se os tradutores investigados podem ser usados como ferramenta de apoio no ensino-aprendizagem e da comunicação de modo geral, no contexto universitário, destacando a qualidade das traduções realizadas pelos sistemas. 

\section{Resultados}

De acordo com \textcite{Santos2018}, as técnicas de IA necessitam de uma grande quantidade de conhecimentos e de mecanismos de manipulação de símbolos para que os computadores se tornem capazes de pensar e tomar decisões. Eles devem ter a possibilidade de representação, modificação e ampliação de suas capacidades de interpretação e de análise de dados. Com a crescente demanda por acessibilidade e inclusão de pessoas surdas e com deficiência auditiva, os tradutores de sinais assumem um papel crucial na comunicação. Para tal, são considerados cinco elementos da língua de sinais: pontos de articulação, configuração de mãos, tipos de movimentos, orientação das mãos e expressões faciais.

\textcite{Lobo2017} complementa afirmando que: 

\begin{quote}
    IA envolve várias etapas ou competências como reconhecer padrões e imagens, entender linguagem aberta escrita e falada, perceber relações e nexos, seguir algoritmos de decisão propostos por especialistas, ser capaz de entender conceitos e não apenas processar dados, adquirir “raciocínios” pela capacidade de integrar novas experiências e, pois, se auto aperfeiçoar (“self learning”), resolvendo problemas, ou realizando tarefas \cite[p. 189]{Lobo2017}
\end{quote}

Com base em afirmações de vários autores que tratam das características da IA e dos tradutores de sinais, foi realizada a análise comparativa dos sete tradutores investigados e identificadas semelhanças, particularidades específicas e limitações relacionadas à precisão e à capacidade de interpretação dos sinais. Os resultados obtidos evidenciam tanto os avanços tecnológicos quanto os desafios ainda existentes para a consolidação de tradutores automáticos eficientes e amplamente acessíveis para a comunidade surda.

A descrição e os resultados foram organizados em duas categorias de tradutores: unidirecional e bidirecional. Na categoria unidirecional, foram agrupados o \textit{Dicionário da Língua Brasileira de Sinais} e o dicionário \textit{Spreadthesign}, os tradutores com uso de luvas \textit{Glove-Based Systems} e os tradutores que usam avatares \textit{Hand Talk} e \textit{VLibras}. Na categoria bidirecional, foram agrupados os tradutores que usam câmeras em tempo real \textit{Sign-Speak} e CESAR + Lenovo. Por fim, foi feita a análise interpretativa dos resultados, apresentando as considerações finais e a conclusão. As comparações foram realizadas com base nas informações contidas nos sites dos respectivos tradutores de sinais investigados e nas pesquisas bibliográficas realizadas.

\subsection{Categoria Unidirecional - Comparação entre o \textit{Dicionário da Língua Brasileira de Sinais} e o \textit{Spreadthesign}}

O \textit{Dicionário da Língua Brasileira de Sinais} e o \textit{Spreadthesign} não são tradutores propriamente ditos. De acordo com informações obtidas nos sites dos referidos dicionários, cada um, possui um banco de dados que funciona como um dicionário visual. Porém, contribuem com igual relevância para a acessibilidade e inclusão comunicacional entre pessoas surdas e ouvintes. Eles são ferramentas digitais voltadas para a aprendizagem e consulta de línguas de sinais, porém com focos e públicos diferentes. Ambos ainda são pouco utilizados em ambientes corporativos, sendo mais comuns entre educadores e estudantes.

As duas são gratuitas, possuem alta acessibilidade e são bastante utilizadas em ambientes educacionais. Suas principais limitações estão no fato de não traduzirem expressões idiomáticas complexas e dependerem da pesquisa por palavras individuais.

No que se refere ao público-alvo, o \textit{Dicionário da Língua Brasileira de Sinais} é mais direcionado para quem busca aprofundamento específico em Libras dentro do contexto brasileiro, enquanto o \textit{Spreadthesign} é recomendado para quem precisa de uma ferramenta mais ampla, com acesso a várias línguas de sinais, em um contexto internacional.

\subsubsection{\textit{Dicionário da Língua Brasileira de Sinais} Unidirecional}

O \textit{Dicionário da Língua Brasileira de Sinais} foi criado por Guilherme de Azambuja Lira e Tanya Amara Felipe de Souza e é mantido pelo INES (Instituto Nacional de Educação de Surdos). Tal dicionário é voltado exclusivamente para a Língua Brasileira de Sinais (Libras), tendo como principal público a comunidade surda brasileira, incluindo estudantes e educadores. Sua interface é em português e adaptada à cultura do Brasil. Oferece vídeos em alta qualidade com sinalizadores surdos nativos, além de definições, frases e categorias temáticas. Tem alta precisão na representação dos sinais e é amplamente utilizado em contextos educacionais no país.

A interface do dicionário apresenta um campo de pesquisa para digitar palavras em português. Ao buscar um termo, o site exibe vídeos demonstrando o sinal correspondente em Libras. Os vídeos são curtos e repetem o movimento do sinal para facilitar a visualização. Além disso, podem aparecer descrições textuais e orientações sobre a posição das mãos e expressões faciais. A navegação é simples, com menus organizados para acessar categorias, verbetes e materiais educativos. Essas características sinalizam que o dicionário pode ser usado como recurso didático desde a educação básica até o ensino superior.

\subsubsection{\textit{Spreadthesign} - Unidirecional}
O dicionário \textit{Spreadthesign} foi criado na Suécia e é mantido pelo European Sign Language Centre. Trata-se de uma plataforma multilíngue que inclui Libras e diversas outras línguas de sinais, atendendo usuários do mundo todo. Está disponível em mais de 40 idiomas e permite a tradução entre línguas faladas e línguas de sinais. Seus vídeos apresentam sinalizadores de diferentes países, o que amplia o repertório de sinais.

O \textit{Spreadthesign} oferece diferentes formas de busca, como por texto, por categorias ou por temas. O usuário pode digitar uma palavra e visualizar imediatamente vídeos com o sinal correspondente. Os vídeos mostram pessoas realizando os sinais de forma clara, com possibilidade de pausar e repetir. Também é possível alternar entre diferentes países para comparar variações da mesma palavra em várias línguas de sinais.

A navegação é intuitiva, com ícones bem definidos e opções de filtro que facilitam a exploração do conteúdo. No entanto, limita-se apenas à tradução da palavra de acordo com o contexto, sem apresentar o significado dela como é feito no Dicionário da Libras. Os referidos dicionários contribuem significativamente para facilitar a comunicação entre pessoas surdas e ouvintes. Além disso, podem ser utilizados como material didático no processo de ensino-aprendizagem, considerando vários aspectos, conforme apresentados na \Cref{tab3}.

\begin{table}[htbp]
\centering
\small
\caption{Contributos dos tradutores unidirecionais do \textit{Dicionário da Língua Brasileira de Sinais} e \textit{Spreadthesign}.}
\label{tab3}
\begin{threeparttable}

\begin{tabularx}{\textwidth}{@{}>{\raggedright\arraybackslash}p{2.5cm}
>{\raggedright\arraybackslash}X
>{\raggedright\arraybackslash}X@{}}

\toprule
Contributo &
Dicionário da Língua Brasileira de Sinais &
Spreadthesign \\
\midrule

Tecnológico &
Plataforma digital nacional baseada em vídeos reais de sinalizantes, acessível via navegador. Estrutura simples e estável, focada no registo lexical. &
Plataforma internacional multimodal (texto + vídeo), com base de dados extensa e suporte a múltiplas línguas de sinais. \\

\addlinespace

Linguístico &
Documenta o léxico da Língua Brasileira de Sinais (Libras), auxiliando na padronização e no ensino do vocabulário. Atua no nível lexical, não frásico. &
Registra e compara sinais de diversas línguas de sinais, promovendo estudos linguísticos contrastivos e evidenciando a diversidade linguística gestual. \\

\addlinespace

Social &
Apoia a aprendizagem de Libras por estudantes, professores, familiares e profissionais, fortalecendo a valorização da comunidade surda brasileira. &
Favorece a comunicação intercultural e internacional de pessoas surdas, promovendo inclusão em contextos globais e acadêmicos. \\

\addlinespace

Aspectos éticos &
Uso ético por apresentar sinais humanos reais. Não substitui intérpretes e deve respeitar variações regionais e contextuais da Libras. &
Valoriza a diversidade cultural e linguística, evitando a padronização excessiva. Também não substitui intérpretes profissionais. \\

\addlinespace

Avanço da comunicação inclusiva &
Amplia o acesso ao conhecimento linguístico em Libras, contribuindo para práticas educacionais inclusivas no contexto nacional. &
Contribui para a inclusão comunicacional em escala global, reconhecendo as línguas de sinais como línguas naturais e distintas. \\

\addlinespace

Custos e acessibilidade econômica &
Gratuito e de livre acesso, adequado para escolas públicas e instituições com poucos recursos. &
Gratuito para consulta online, com ampla acessibilidade econômica e institucional. \\

\addlinespace

Perspectivas futuras &
Expansão do vocabulário, maior cobertura de variações regionais e integração com plataformas educacionais digitais. &
Inclusão de mais línguas de sinais, fortalecimento da pesquisa linguística e possível integração com tecnologias de IA educacional. \\

\addlinespace

Integração &
Integração Fácil integração em cursos de Libras, materiais didáticos e ambientes virtuais de aprendizagem. &
Integração em contextos acadêmicos internacionais, formação de intérpretes e projetos de cooperação linguística. \\

\bottomrule

\end{tabularx}
\source{Elaboração própria, com base nos sites dos dois tradutores investigados e nas pesquisas bibliográficas realizadas.}
\end{threeparttable}
\end{table}

\subsection{Tradutor unidirecional – \textit{Glove-Based Systems}}

O projeto \textit{Glove-Based Systems} foi desenvolvido por pesquisadores da Universidade da Califórnia em Los Angeles (UCLA) e é mantido pela própria instituição. Ele possui características distintas dos demais tradutores, embora siga os mesmos critérios de análise considerando a inclusão comunicacional da pessoa surda.

Os \textit{Glove-Based Systems} são tecnologias experimentais que utilizam luvas com sensores para reconhecer e traduzir sinais manuais da língua de sinais para a língua oral ou escrita, sem realizar o processo inverso. Esses sistemas registram com precisão os movimentos dos dedos e a posição das mãos, mas ainda têm baixa a moderada precisão linguística, pois não conseguem representar toda a complexidade das línguas de sinais, sobretudo as expressões faciais e corporais.

De acordo com \textcite{LuEtAl2023}, a luva inteligente vestível de língua de sinais fornece uma nova estratégia para as pessoas com deficiência auditiva se comunicarem com o mundo. As luvas são portáteis e funcionam em diferentes ambientes, sem depender de iluminação de fundo, mas enfrentam desafios como desconforto no uso prolongado, peso, limitação de movimentos natural das mãos, alto custo, necessidade de calibração individual e pouca acessibilidade em países em desenvolvimento. Projetos de baixo custo ou de código aberto ainda são raros.

\textcite{LiangJettanasenChiradeja2024} destaca a importância da padronização dos sinais para facilitar a comunicação entre pessoas com deficiência auditiva, embora ainda existam desafios nesse processo. Segundo o autor, muitas pessoas da comunidade surda e com deficiência auditiva enfrentam dificuldades para compreender a língua de sinais porque não dominam suficientemente o sistema de sinais, ou seja, não possuem conhecimento adequado da forma correta de sinalizar (vocabulário, estrutura e uso da língua). Por isso, o seu público-alvo inclui pessoas surdas, instituições educacionais, organizações e pesquisadores.

Atualmente, esses sistemas não são multilíngues, pois cada língua de sinais possui estrutura, gestos e gramática próprios, exigindo o treinamento específico dos algoritmos. Embora não sejam projetados originalmente para reconhecer expressões faciais e corporais, pesquisadores já realizam testes com sensores adesivos posicionados no rosto para captar parte dessas expressões, que são fundamentais na comunicação em língua de sinais.

A análise também denotou que as luvas tradutoras, apesar das limitações, podem ser úteis em contextos específicos, como na educação, facilitando a comunicação entre crianças surdas e professores ouvintes; em serviços públicos e atendimentos rápidos sem intérpretes; e em ambientes industriais com muito ruído. Entretanto, o impacto social é desigual, pois a tecnologia tende a priorizar a adaptação da pessoa surda ao mundo ouvinte, em vez de valorizar o bilinguismo e a língua de sinais como língua legítima.

Outra questão percebida na análise foi que, apesar da inovação, o contributo linguístico dos sistemas, baseados em luvas enfrentam desafios críticos quando confrontados com a complexidade linguística da LS, como a redução do sinal linguístico, que envolve, também, expressões faciais, movimentação corporal, orientação do tronco e ritmo. Nesses aspectos, as luvas não conseguem captar sozinhas, desencadeando uma tradução incompleta ou imprecisa e perdendo nuances semânticas e gramaticais (como interrogativas marcadas pelo rosto ou tempos verbais marcados pelo espaço).

Do ponto de vista ético e da sustentabilidade, percebeu-se que há preocupações com a dependência de dispositivos corporais, que podem causar estigmatização, e com a visão capacitista, que prioriza a tradução da LS para ouvintes em vez de incentivar o aprendizado da própria língua de sinais. Do ponto de vista da ergonomia e da usabilidade, muitos modelos são pesados, desconfortáveis e restringem os movimentos naturais das mãos, causando fadiga e dificultando a adesão dos usuários. Além disso, essas luvas apresentam limitações práticas: alto custo, baixa acessibilidade em países em desenvolvimento, pouca disponibilidade de suporte técnico e escassez de projetos de baixo custo e código aberto.

A evolução futura dos sistemas baseados em luvas dependerá da integração com tecnologias complementares, como a visão computacional, para captar expressões faciais; redes neurais multimodais, que unificam sinais manuais e não manuais; interfaces neurais e hápticas, que oferecem retorno tátil adaptado ao usuário; e realidade aumentada ou hologramas, para melhorar a interpretação dos sinais e a experiência do usuário.

\subsection{Comparação dos tradutores unidirecionais –\textit{Vlibras} e \textit{Hand Talk}}
O \textit{VLibras }e o \textit{Hand Talk} são ferramentas brasileiras desenvolvidas para tornar a comunicação em Libras mais acessível e possuem propostas e formatos de uso diferentes. Ambas as plataformas contam com suporte à língua de sinais e avatares animados, mas enfrentam desafios relacionados à precisão, especialmente na tradução de expressões idiomáticas e contextos complexos, pois as línguas de sinais envolvem elementos visuais e culturais que nem sempre são totalmente representados por avatares. Elas, também não captam expressões faciais ou corporais para tradução, trabalhando apenas com entrada por texto ou voz.

De modo geral, nota-se que o VLibras se destaca pela gratuidade, caráter público e integração governamental, voltada à acessibilidade digital no Brasil, enquanto o Hand Talk se consolida como uma solução comercial mais expressiva visualmente e com alcance internacional, incluindo suporte à ASL.

\subsubsection{Tradutor unidirecional – \textit{VLibras}}
O \textit{VLibras }é um projeto público criado pela Universidade Federal da Paraíba (UFPB) em parceria com o Ministério da Gestão e Inovação e outros órgãos federais. Seu principal objetivo é tornar conteúdos digitais acessíveis em Libras, especialmente em sites institucionais e governamentais. Ele utiliza processamento de linguagem natural e avatar 3D para traduzir textos e áudios para Libras. Possui três avatares (Ícaro, Hozana e Guga) e oferece interação quase em tempo real.

Está disponível gratuitamente para web, desktop e dispositivos móveis. Sua precisão pode variar, especialmente em expressões idiomáticas e contextos complexos. Embora não seja multilíngue, permite adaptar as traduções a regionalismos de diferentes estados brasileiros. O sistema não capta expressões faciais ou corporais do usuário; a tradução ocorre apenas a partir de entrada textual ou de voz.

\subsubsection{Tradutor unidirecional – \textit{Hand Talk}}
O Hand Talk, foi criado pela startup brasileira Hand Talk Tecnologia S.A. e tem como objetivo facilitar a comunicação entre surdos e ouvintes, especialmente, para empresas e escolas, mas, também pode ser usado pelo público de modo geral. Ele oferece aplicativos móveis e Interface de Programação de Aplicações (API) comercial para realizar traduções de texto e voz e utiliza inteligência artificial combinada a avatar 3D. Possui dois avatares adultos (Hugo e Maya) animados que oferecem suporte a duas línguas de sinais: Libras e ASL (Língua de Sinais Americana). 

Sua tradução não ocorre plenamente em tempo real, mas é eficiente para diversos contextos de atendimento e acessibilidade digital. Esses dois tradutores podem contribuir para a comunicação entre pessoas surdas e ouvintes. Além disso, podem ser utilizados como material didático no processo de ensino-aprendizagem, conforme apresentado na \Cref{tab4}.

\begin{table}[htbp]
\centering
\small
\caption{Contributos dos tradutores unidirecional Hand Talk e VLibras.}
\label{tab4}
\begin{threeparttable}

\begin{tabularx}{\textwidth}{@{}>{\raggedright\arraybackslash}p{2.8cm}
>{\raggedright\arraybackslash}X
>{\raggedright\arraybackslash}X@{}}

\toprule
Contributo &
Hand Talk &
VLibras \\
\midrule

Tecnológico &
Plataforma privada com avatar 3D (Hugo/Maya), \textit{app} móvel, \textit{plugin} para sites e API para integração empresarial. Foco em experiência visual e produtos comerciais. &
Conjunto de ferramentas \textit{open-source} (\textit{plugin} web, \textit{desktop}, API). Desenvolvido com foco em serviços públicos e integração ampla em portais governamentais. \\

\addlinespace

Linguístico &
Tradução automática do português para Libras (e outras línguas de sinais). Adequada para frases simples e vocabulário básico; limitada em construções gramaticais complexas. &
Tradução automática do português para Libras com foco linguístico nacional. Melhor adaptação à Libras, mas ainda com limitações sintáticas e semânticas em textos complexos. \\

\addlinespace

Social &
Populariza o contacto com Libras entre ouvintes, escolas e empresas. Aumenta a sensibilização social para a inclusão de pessoas surdas. &
Amplia o acesso da comunidade surda a informações públicas, educação e serviços governamentais, promovendo inclusão digital em larga escala. \\

\addlinespace

Aspectos éticos &
Risco de uso inadequado como substituto a intérpretes humanos. Traduções automáticas podem gerar erros de sentido. Questões de privacidade dependendo da integração em nuvem. &
Ética baseada no acesso público e gratuito. Também não substitui intérpretes humanos; exige uso responsável, sobretudo em contextos críticos (saúde, justiça, educação formal). \\

\addlinespace

Avanço da comunicação inclusiva &
Facilita a comunicação básica entre surdos e ouvintes, especialmente em ambientes digitais privados e educativos. Forte impacto em acessibilidade corporativa. &
Promove comunicação inclusiva em políticas públicas, educação e cidadania digital, tornando a Libras visível em sites oficiais. \\

\addlinespace

Custos e acessibilidade econômica &
Não é totalmente gratuito, tendo em vista que possui versão gratuita limitada, sendo sua \textit{Application Programming Interface} (API) e recursos mais avançados (soluções empresariais) pagos. &
Totalmente gratuito e de código aberto, reduzindo barreiras econômicas para instituições públicas e organizações sociais. \\

\addlinespace

Perspectivas Futuras &
Evolução com IA mais avançada, avatares mais naturais, expansão internacional e integração multimodal (voz, texto, vídeo). &
Crescimento através da comunidade \textit{open-source}, melhoria linguística contínua e maior integração com plataformas educacionais e governamentais. \\

\addlinespace

Integração &
Forte integração com websites, \textit{apps}, sistemas corporativos e produtos educacionais por meio de APIs e SDKs. &
Integração ampla em portais públicos, universidades, plataformas governamentais e sistemas educacionais, favorecendo a interoperabilidade. \\

\bottomrule

\end{tabularx}
\source{Elaboração própria com base nos sites dos dois tradutores investigados e nas pesquisas bibliográficas realizadas.}
\end{threeparttable}
\end{table}

\subsection{Categoria Bidirecional – Comparação dos tradutores \textit{Sign-Speak} e Projeto CESAR + Lenovo}
Na comparação dos tradutores Sign-Speak e Projeto CESAR +Lenovo, percebe-se que ambos se destacam porque usam como tecnologia principal câmeras do tipo RGB/3D ou similar, com visão computacional e rastreamento corporal. A depender da resolução e do ângulo da câmera, a captação de sinais manuais é boa. A captação de expressões faciais e corporais é alta, pois permite leitura facial e espacial (não-manual). A fidelidade linguística pode ser considerada alta devido à boa captação de elementos gramaticais da LS. Subtende-se que esses tradutores possam ser usados com sucesso, na educação e nas organizações de atendimento ao público.

Esses dois tradutores, \textit{Sign-Speak} e Projeto CESAR + Lenovo, são projetos que oferecem interação em tempo real, mas ainda enfrentam desafios técnicos relacionados à precisão, latência e às variações de iluminação, fundo, ângulo da câmera e vestuário. A precisão atual é considerada adequada, principalmente para comunicações simples e informais. De modo geral, são tecnologias promissoras, mas que ainda estão em processo de amadurecimento antes de alcançarem ampla adoção no cotidiano.

\subsubsection{Tradutor bidirecional – \textit{Sign-Speak}}
O Sign-Speak é um conjunto de projetos internacionais, desenvolvidos por equipes na Europa e nos Estados Unidos, com o objetivo de permitir a comunicação automática entre pessoas surdas e ouvintes. Ele utiliza reconhecimento facial e de sinais por meio de câmeras para traduzir sinais para texto e voz, e vice-versa. Seu foco principal são as línguas de sinais britânica e americana.

O sistema é projetado para reconhecer movimentos das mãos, expressões corporais e aspectos visuais da comunicação. Não utiliza avatares animados, pois a comunicação é baseada diretamente na captura de imagens das pessoas. O \textit{Sign-Speak} não funciona como dicionário, mas sim como um sistema de interpretação em tempo real. Atualmente, está em fase de desenvolvimento e testes, com alguns protótipos disponíveis, sendo voltado tanto para usuários surdos e ouvintes quanto para desenvolvedores de soluções inclusivas.

\subsubsection{Tradutor bidirecional – Projeto CESAR + Lenovo}

O projeto CESAR + Lenovo é uma iniciativa brasileira desenvolvida em parceria entre o CESAR (Centro de Estudos e Sistemas Avançados) e a empresa Lenovo. Seu foco principal é a tradução da Libras para texto e fala em tempo real, visando promover a inclusão, facilitar a comunicação em ambientes sem intérpretes e apoiar o ensino da língua de sinais. O sistema funciona por meio da captura dos sinais feitos pela pessoa, analisando movimentos das mãos, articulações, velocidade e contexto do sinal para gerar frases completas em português. Assim como o \textit{Sign-Speak}, não utiliza avatares animados, pois se baseia diretamente na imagem do usuário captada pela câmera.

O projeto CESAR + Lenovo está em fase de testes (prova de conceito) e ainda não há confirmação de sua disponibilização gratuita ao público. Sua proposta é atender principalmente pessoas surdas em contextos de atendimento, educação e serviços, podendo ser utilizado em empresas, escolas e instituições públicas. Diferente de um dicionário, ele busca interpretar frases completas de forma contextual e não apenas traduzir sinais isolados.

Os tradutores bidirecionais \textit{Sign-Speak} e CESAR + Lenovo contribuem em vários aspectos para o avanço comunicacional inclusivo tecnológico, linguístico, social, ético, entre outros, conforme apresentado na \Cref{tab5}.

\begin{table}[htbp]
\centering
\small
\caption{Contributos dos tradutores bidirecional Sign-Speak e CESAR + Lenovo.}
\label{tab5}
\begin{tabularx}{\textwidth}{>{\raggedright\arraybackslash}p{0.20\textwidth}
                              >{\raggedright\arraybackslash}X}
\toprule
Contributo & Tradutores bidirecionais – Sign-Speak e CESAR + Lenovo \\
\midrule

Tecnológico &
As tecnologias utilizam câmeras 3D para possibilitar o rastreamento corporal e captar gestos em tempo real, promovendo a inclusão comunicacional entre surdos e ouvintes por meio da tradução bidirecional com sinais e fala ou texto. A conectividade sem fio permite a comunicação com dispositivos móveis, o que também facilita a comunicação em locais variados.
\\[0.5em]

Linguístico &
Ambos ajudam a destacar a complexidade linguística da LS, combatendo visões preconceituosas que a consideram “inferior” ou “incompleta”. Além disso, ele pode contribuir com a documentação linguística de variações e dialetos da LS e o desenvolvimento de corpus digital útil para linguistas, educadores e desenvolvedores, tendo em vista que a LS não é apenas uma versão gestual das línguas faladas.
\\[0.5em]

Social &
O principal contributo dos tradutores Sign-Speak e CESAR + Lenovo é promover a inclusão digital e social de pessoas surdas, permitindo a comunicação em tempo real em diversas situações do dia a dia. A ferramenta favorece a autonomia comunicacional, reduzindo a dependência de intérpretes, facilita a integração de estudantes surdos na educação regular com traduções automáticas e amplia o acesso ao mercado de trabalho em ambientes com predominância da comunicação verbal.
\\[0.5em]

Ético &
Do ponto de vista ético, é importante considerar que, apesar dos benefícios, tecnologias como tradutores automáticos enfrentam desafios como: riscos à privacidade devido à captação constante de dados; necessidade de representatividade nos modelos, incluindo a diversidade da comunidade surda; risco de dependência excessiva da tecnologia em contextos sensíveis; e o uso superficial da inclusão, sem promover mudanças reais nas estruturas sociais excludentes.
\\[0.5em]

Avanços na comunicação inclusiva &
Os tradutores Sign-Speak e CESAR + Lenovo não devem ser vistos apenas como ferramentas de tradução, mas como o início de uma revolução comunicacional inclusiva. Eles ampliam o reconhecimento da língua de sinais como legítima e essencial na construção de uma sociedade mais equitativa. Ao facilitar o diálogo entre surdos e ouvintes, contribuem para reduzir barreiras históricas de exclusão. Além disso, estimulam o desenvolvimento de tecnologias mais sensíveis às diversidades linguísticas e culturais da comunidade surda.
\\[0.5em]

Custos e acessibilidade econômica &
O contributo na perspectiva dos custos e da acessibilidade econômica deve ser amplo e viável, especialmente em contextos de desigualdade. Tecnologias assistivas ainda são elitizadas, perpetuando exclusões. A inclusão via inteligência artificial requer políticas públicas, apoio privado e investimentos que garantam acesso universal, com foco na equidade e na participação ativa da comunidade surda, respeitando sua cultura e identidade.
\\[0.5em]

Perspectivas futuras de integração &
O seu contributo, em perspectivas futuras de integração, é o de que a inteligência artificial pode tornar a sociedade mais inclusiva ao oferecer acesso personalizado e integrar grupos marginalizados. Para isso, é preciso alinhar o desenvolvimento da IA a princípios éticos, à participação ativa dos beneficiados e ao foco em necessidades reais. Expressões como “perspectivas futuras de integração” devem refletir tanto o papel social da tecnologia quanto a inclusão plena, respeitando a diversidade cultural, linguística e funcional.
\\

\bottomrule
\end{tabularx}
\source{Elaboração própria com base nos sites dos dois tradutores investigados e nas pesquisas bibliográficas realizadas.}
\end{table}

A \Cref{tab6} permite observar que a evolução dos tradutores depende essencialmente da capacidade das tecnologias futuras de integrar movimento, configuração de mão, orientação espacial, expressão facial e contexto semântico (significado das palavras, frases e expressões, analisando como elas transmitem o sentido e como o contexto afeta essa interpretação). Esses elementos, quando combinados, definem a complexidade única das línguas de sinais.

% {\def\LTcaptype{none} % do not increment counter
% \begin{longtable}[]{@{}
%   >{\raggedright\arraybackslash}p{(\linewidth - 8\tabcolsep) * \real{0.1999}}
%   >{\centering\arraybackslash}p{(\linewidth - 8\tabcolsep) * \real{0.1971}}
%   >{\centering\arraybackslash}p{(\linewidth - 8\tabcolsep) * \real{0.2028}}
%   >{\centering\arraybackslash}p{(\linewidth - 8\tabcolsep) * \real{0.1946}}
%   >{\raggedright\arraybackslash}p{(\linewidth - 8\tabcolsep) * \real{0.2054}}@{}}
% \toprule\noalign{}
% \endhead
% \bottomrule\noalign{}
% \endlastfoot
\begin{table}[htbp]
\centering
\footnotesize
\caption{Tradutores de línguas de sinais: comparativo transversal.}
\label{tab6}
\begin{threeparttable}
\begin{tabularx}{\textwidth}{
    >{\raggedright\arraybackslash}p{0.2\textwidth}
    >{\centering\arraybackslash}p{0.15\textwidth}
    >{\centering\arraybackslash}p{0.15\textwidth}
    >{\centering\arraybackslash}p{0.15\textwidth}
    >{\raggedright\arraybackslash}p{0.2\textwidth}
}

\toprule
Contributo &
Hand Talk &
VLibras \\
\midrule
\multirow{2}{=}{Critério transversal} &
\multicolumn{3}{>{\centering\arraybackslash}p{0.45\textwidth}}{Tradutores unidirecionais} & Tradutores bidirecionais \\
& \emph{Dicionários}
\emph{Libras} e \emph{Spreadthesign} & Luvas/sensores & Automáticos (texto - avatar) & IA (visão-fala/texto) \\
\midrule
Gratuitos/pagos & Maioritariamente gratuitos & Variáveis
(muitos requerem compra de hardware) & Maioritariamente gratuitos &
Variáveis (alguns gratuitos, outros experimentais) \\
Necessidade de hardware especial? & Não & Sim (luvas/ sensores)
& Não & Não (alguns usam apenas webcam) \\
Precisão funcional geral & Média (bons para palavras isoladas)
& Baixa (limitações extremas de reconhecimento) & Média (tradução
limitada e artificial) & Variável (dependem do algoritmo; ainda
instáveis) \\
Adequado para frases completas? & Não & Não & Não &
Parcialmente (com muita limitação) \\
Captura expressões faciais/não manuais? & Não & Não & Não &
Parcialmente (ainda fracos) \\
Facilidade de uso & Alta & Baixa & Média & Média (interface
simples) \\
Acessível para contexto escolar? & Sim & Não & Sim &
Parcialmente (dependem da internet e do desempenho) \\
Funcionamento offline? & Parcial (alguns) & Não & Não & Não \\
Aplicação educativa prática & Muito úteis & Muito limitados &
Úteis para sensibilização & Promissores, mas não fiáveis \\
Velocidade (em tempo real?) & Não & Teoricamente sim, porém
imprecisos & Geram avatares em tempo quase real & Sim, com latência \\
Tipo de tradução & Palavras isoladas & Gestos isolados & Texto
- avatar & Gesto - texto/fala; texto - gesto \\
\bottomrule
\end{tabularx}
\source{Elaboração própria com base nos sites dos dois tradutores investigados e nas pesquisas bibliográficas realizadas.}
\end{threeparttable}
\end{table}
%\end{longtable}
%}


O resultado do comparativo transversal dos tradutores evidencia que cada categoria de tradutor opera com pressupostos tecnológicos diferentes, o que limita comparações diretas, mas permite identificar tendências, padrões de funcionamento e fragilidades comuns.

\section{Análise interpretativa dos resultados}
De acordo com \textcite{Ahmed2018}, a língua de sinais (LS) é visual e espacial, formada por componentes posicionais e visuais, como o formato dos dedos e das mãos, a localização e orientação das mãos, além dos movimentos dos braços e do corpo. Esses elementos atuam juntos para transmitir significados. Sua estrutura fonológica geralmente inclui cinco elementos fundamentais, que servem como blocos de construção de cada sinal. Esses blocos podem ser utilizados por sistemas inteligentes para realizar o reconhecimento automático da língua de sinais \textit{(Sign Language Recognition} – SLR).

Foi com base nessa premissa que foi feita a análise interpretativa dos resultados obtidos na investigação. A análise também considerou a diversidade tecnológica observada nos diferentes tipos de tradutores da língua de sinais - dicionários visuais, dispositivos com sensores (luvas), tradutores unidirecionais automáticos e sistemas bidirecionais baseados em visão computacional e revela que essa área se encontra num processo de desenvolvimento assimétrico e sem um padrão uniforme de qualidade.

Os dicionários digitais tradutores da língua de sinais continuam a ser a ferramenta mais consistente e pedagogicamente funcional. Embora, raramente, permitam a tradução de frases completas, apresentam grande facilidade de uso, acessibilidade e precisão razoável para palavras isoladas, o que os tornam adequados para o contexto escolar. Esses recursos têm a vantagem de serem gratuitos e estáveis, mas não alcançam níveis de tradução contextual ou expressiva.

Os dispositivos com luvas ou sensores revelam-se a categoria com pior desempenho transversal. Apesar do potencial teórico de traduzir movimento de forma direta, a tecnologia atual demonstra baixa precisão, dependência absoluta de hardware especializado e ausência de reconhecimento facial ou corporal, aspetos essenciais para a tradução fiel de língua de sinais. Como resultado, esta categoria não apresenta viabilidade prática, seja em contexto educativo, comunicacional ou profissional.

Os tradutores unidirecionais automáticos (texto - avatar) apresentam vantagens em termos de acessibilidade e simplicidade de uso, permitindo converter frases escritas em animações por meio de uso de sinais. Contudo, estes sofrem de rigidez linguística, pouca naturalidade dos movimentos e ausência de expressões faciais, o que resulta numa tradução artificial e, em alguns casos, potencialmente enganadora. São úteis para a sensibilização e a exploração inicial, mas não substituem intérpretes nem suportam comunicação real.

A categoria bidirecional, a mais avançada e investigada atualmente, apresenta o maior potencial, sobretudo por integrar visão computacional e aprendizagem automática. No entanto, os resultados mostram que essa tecnologia ainda carece de maturidade: embora consiga captar alguns movimentos e produzir tradução parcial para texto ou voz, a precisão é altamente variável, especialmente devido à dificuldade em reconhecer simultaneamente mãos, configuração dos dedos, movimento corporal e expressões faciais. A promessa tecnológica é elevada, mas a fiabilidade ainda não permite uso autônomo em ambientes profissionais ou educativos exigentes. 

De forma transversal, os resultados revelam quatro conclusões principais:

\begin{enumerate}
    \item nenhuma tecnologia atual consegue realizar tradução completa, natural e contextual da língua de sinais;
    \item as categorias mais avançadas tecnologicamente são também as menos estáveis, devido à complexidade do reconhecimento gestual integral;
    \item a acessibilidade escolar favorece dicionários e tradutores unidirecionais, enquanto luvas e sistemas bidirecionais apresentam maiores barreiras técnicas;
    \item as expressões faciais, essenciais na gramática da língua de sinais, permanecem um desafio crítico, sobretudo fora dos sistemas bidirecionais.
    \item 
\end{enumerate}

A solução mais funcional hoje continua a ser a combinação de recursos digitais complementares, não a substituição total por um único tradutor automático. Desta forma, a análise interpretativa evidencia que, apesar de avanços consideráveis, a tradução automática de língua de sinais permanece tecnologicamente fragmentada e inconstante. As ferramentas disponíveis hoje servem sobretudo para apoio pedagógico, consulta lexical, experimentação tecnológica e sensibilização, mas não garantem ainda equivalência comunicativa nem precisão linguística.

\section{Conclusão}\label{sec-conclusao}
Em conclusão, os objetivos propostos foram plenamente atendidos, uma vez que foi possível analisar e comparar sete tradutores de Língua de Sinais: \textit{Dicionário da Língua Brasileira de Sinais},  \textit{Spreadthesign, Glove-Based Systems, Hand Talk, VLibras, Sign-Speak} e CESAR + Lenovo, com base em três eixos analíticos: (1) recursos de dicionário (tradutor automático); (2) tradução automática unidirecional; e (3) tradução automática bidirecional, bem como a sua aplicabilidade em contextos institucionais, educacionais e corporativos.

A análise evidenciou que, apesar dos avanços tecnológicos e do crescente interesse científico na área, a tradução automática da língua de sinais ainda apresenta limitações significativas em termos de precisão, naturalidade e fiabilidade. Dessa forma, conclui-se que essas ferramentas, embora úteis como apoio, ainda não são capazes de substituir intérpretes humanos nem de garantir uma comunicação plenamente autônoma.

Os dicionários digitais permanecem as ferramentas mais estáveis e funcionalmente úteis no contexto educativo enquanto tecnologias mais avançadas, como sistemas bidirecionais baseados em visão computacional, revelam grande potencial, mas exibem ainda limitações significativas em termos de reconhecimento de expressões faciais, captação do movimento global e interpretação semântica, aliadas ao fato de que a gramática das línguas de sinais é diferente da língua oral e os erros de interpretação ainda são comuns.

Apesar das limitações ainda presentes nos tradutores de língua de sinais online com recurso à inteligência artificial, como dificuldades na interpretação de nuances linguísticas, expressões idiomáticas e contextos acadêmicos complexos, essas tecnologias vêm desempenhando um papel cada vez mais relevante na facilitação da comunicação entre pessoas surdas e ouvintes.

De acordo com \textcite{PapastratisEtAl2021}, as tecnologias de inteligência artificial voltadas para o reconhecimento e a tradução da língua de sinais têm um contributo significativo na redução das barreiras comunicacionais, ampliando o acesso à informação e promovendo a inclusão social de pessoas surdas.

Embora esses tradutores sejam importantes para a interação entre surdos e ouvintes, devem ser usados, no contexto educativo, apenas como ferramentas complementares. São úteis para facilitar o acesso a conteúdos acadêmicos, a plataformas online, materiais didáticos online e interações e mediação pedagógica em ambientes virtuais de aprendizagem, promovendo maior autonomia dos estudantes surdos, para explorar o vocabulário e promover a sensibilização para a língua de sinais, mas não substituem a comunicação formal. Sua utilização deve ser criteriosa, acompanhada de reflexão crítica sobre suas limitações e sempre integrada a práticas pedagógicas inclusivas, com a participação ativa de usuários fluentes.

\section{Recomendações para pesquisas futuras}
Para a investigação futura, tornam-se essenciais abordagens que articulem tecnologia, linguística e educação, explorando novas formas de reconhecimento multicanal de sinais, análises corpóreas mais profundas e integração de expressões não manuais. A investigação deve também recorrer a metodologias participativas, envolvendo pessoas surdas nos processos de conceção e validação dos sistemas, garantindo que as soluções tecnológicas reflitam a complexidade real da língua de sinais. Só através desta articulação multidisciplinar será possível evoluir para ferramentas mais robustas, éticas e verdadeiramente inclusivas.

\printbibliography\label{sec-bib}
% if the text is not in Portuguese, it might be necessary to use the code below instead to print the correct ABNT abbreviations [s.n.], [s.l.]
%\begin{portuguese}
%\printbibliography[title={Bibliography}]
%\end{portuguese}


%full list: conceptualization,datacuration,formalanalysis,funding,investigation,methodology,projadm,resources,software,supervision,validation,visualization,writing,review
\begin{contributors}[sec-contributors]
\authorcontribution{Maria Aparecida Araujo Lima}[conceptualization,datacuration,investigation,validation,resources,writing,review]
\authorcontribution{Ana Garcia-Valcarcel
Muñoz-Repiso}[validation,resources,review]
\authorcontribution{Manuel Meirinhos}[projadm,validation,resources,review]
\end{contributors}

\begin{dataavailability}
\txtdataavailability{databody} % options: dataavailable, dataonly, databody, datanotav, nodata
\end{dataavailability}

\section*{\fontsize{10}{12}\selectfont Agradecimentos}
Esta pesquisa foi apoiada pela Universidade de Salamanca, pela Universidade Politécnica de Bragança e pela Universidade Federal de Alagoas (UFAL).

\end{document}

